\section{Implementation}
\label{sec:implementation}

We implemented \tool{} as a modular analysis pipeline. The current prototype
contains four components: a standard parser, a code indexer, an LLM orchestration
layer, and a validation-test runner. This section describes implementation
choices that affect reproducibility and engineering trade-offs.

\textbf{Standard processing.}
The standard parser ingests RFC-style text, markdown, or converted PDF text. It
segments documents by section, preserves sentence offsets, and detects
normative keywords using a conservative lexicon derived from RFC 2119-style
language. Field candidates are extracted from message diagrams, bullet lists,
tables, and definitional paragraphs. When a field candidate is ambiguous, the
parser retains all candidates and asks the extraction stage to resolve the
reference against local context.

\textbf{Code indexing.}
The code indexer supports C and C++ in the prototype. It uses Tree-sitter with
C/C++ grammars for a uniform symbol index and compiler-aware parsing where
available for richer scope information. Each indexed symbol records its file,
scope, declaration span, use spans, surrounding comments, and enclosing
function or type. The indexer also extracts a lightweight call graph from direct
calls and function-pointer-like patterns when they are syntactically visible.

\textbf{LLM orchestration.}
LLM calls are isolated behind typed prompts with JSON-like output schemas. The
pipeline uses low-temperature decoding for extraction, alignment, and
reasoning. Every LLM output passes through a deterministic validator that checks
schema validity, source-span existence, identifier presence, duplicate entries,
and label membership. Invalid outputs are retried with a compact repair prompt;
after the retry budget is exhausted, \tool{} falls back to deterministic
rankings or marks the sample as insufficient evidence.

\textbf{Prompt boundaries.}
The prompts are intentionally narrow. The extractor sees a section and the
field space, not the whole standard. The aligner sees a rule and ranked code
snippets, not the whole repository. The reasoner sees the evidence bundle and
must cite snippet identifiers rather than free-form code claims. This design
keeps token cost manageable and makes the output auditable.

\textbf{Validation runner.}
The validation runner currently supports three modes: extending an existing
unit test, creating a small harness around a parser or validation function, and
driving a command-line implementation with generated inputs. The runner records
compiler options, environment variables, command lines, exit status, stdout,
stderr, and relevant logs. If the target cannot be built or the path cannot be
reached, the sample is not counted as confirmed.

\textbf{Parameters.}
Unless otherwise stated, the prototype keeps the top 8 files after file
ranking, the top 20 variable candidates per rule, and the top 12 snippets per
reasoning round. The maximum number of evidence-completion rounds is three. The
validation agent receives at most two attempts per hypothesis before reporting
unmet validation conditions. We tune these parameters on held-out protocol
sections and keep them fixed for all evaluated implementations.

\textbf{Caching and determinism.}
To make repeated runs auditable, \tool{} stores every prompt, response,
validator decision, and retrieved code bundle under a content hash. When the
same rule and repository snapshot are analyzed again, deterministic stages are
reused and LLM stages can be replayed from cache. The cache is also useful for
manual inspection: a reviewer can open a report and see the exact evidence
available to the reasoner at each round.

\textbf{Failure handling.}
The implementation distinguishes recoverable and terminal failures. Recoverable
failures include malformed model output, missing optional build metadata, or a
snippet that cannot be parsed by one parser but can still be indexed by
Tree-sitter. Terminal failures include missing source files, standards that
cannot be segmented, or implementations that cannot expose any parse or
validation entry point. Terminal failures are reported as coverage gaps rather
than silently omitted from aggregate metrics.

\textbf{Manual review interface.}
Although \tool{} is designed for automation, the prototype emits a compact
review bundle for each candidate. The bundle contains the source text, the
normalized rule, ranked variables, selected snippets, reasoning transcript,
validation commands, and final label. This interface is important for building
ground truth: protocol experts and maintainers can correct a rule, mark an
implementation-specific exception, or confirm that a test demonstrates a real
deviation.
